\documentclass[10pt, handout]{beamer}
\usepackage{graphicx} % Required for inserting images
\usepackage{amsmath, amsthm, amssymb}
\usepackage{xcolor}

% Theme and Color Configuration
% \usetheme{Madrid}
% \usecolortheme{beaver}
% \setbeamertemplate{frame footer}{\insertshortauthor}
% \setbeamertemplate{footline}{\insertshortauthor~(\insertshortinstitute)\hfill\insertshorttitle}

\setbeamerfont{title}{size=\large}
\setbeamerfont{author}{size=\small}
\setbeamerfont{date}{size=\scriptsize}
\setbeamerfont{normal text}{size=\small}
\setbeamerfont{itemize/enumerate body}{size=\small}
\setbeamerfont{itemize/enumerate subbody}{size=\footnotesize}
\setbeamerfont{frametitle}{size=\large}
\setbeamerfont{block title}{size=\normalsize}
\setbeamerfont{block body}{size=\small}
\setbeamerfont{author in head/foot}{size=\tiny}

\usefonttheme{serif}
% \usepackage[scaled=0.95]{newpxtext}
% \usepackage{newpxmath}
\usepackage{mathpazo}
\newcommand{\norm}[1]{\left\lVert#1\right\rVert}
\DeclareMathOperator{\Var}{Var}

\definecolor{darkred}{rgb}{0.55, 0.0, 0.0}
\definecolor{maroon(html/css)}{rgb}{0.5, 0.0, 0.0}
\setbeamercolor{title}{fg=darkred}
\setbeamercolor{frametitle}{fg=darkred}

\setbeamercolor{theorem title}{bg=maroon(html/css)!90!white, fg=white}
\setbeamercolor{theorem body}{bg=maroon(html/css)!20!white, fg=black}
\setbeamercolor{block title}{bg=maroon(html/css)!90!white, fg=white}
\setbeamercolor{block body}{bg=maroon(html/css)!20!white, fg=black}
\setbeamercolor{author in head/foot}{fg=black!30!white}
\setbeamercolor{title in head/foot}{fg=black!30!white}

\setbeamertemplate{itemize item}[ball]
\setbeamercolor{item}{fg=darkred}
\setbeamercolor{itemize subitem}{fg=darkred}
\setbeamercolor{enumerate item}{fg=darkred}
\setbeamercolor{navigation symbols}{fg=darkred, bg=darkred!20!white}
\setbeamertemplate{footline}{%
  \begin{beamercolorbox}[ht=2.5ex,dp=2.5ex,%
    leftskip=.3cm,rightskip=.3cm plus1fil]{title in head/foot}
    \leavevmode{\usebeamerfont{title in head/foot}%
      [\insertshorttitle, \insertframenumber]}
  \end{beamercolorbox}%
}
\setbeamertemplate{navigation symbols}{}

\setbeamertemplate{blocks}[rounded][shadow=true]

\title[Causal Lab]{Causal Inference}
\subtitle[]{Lab 7: Bias-Corrected Matching}
\author[Song and Park]{Xiangyu Song and Chanhyuk Park}
\institute{Washington University in St. Louis}
\date{\today}

\begin{document}

\maketitle

\begin{frame}
  \frametitle{Matching as a Pre-Processing Step}

  \begin{itemize}
    \item We have covered matching in randomized experiments.
    \item Matching can be used as a pre-processing step for causal inference with observational data
    \begin{enumerate}
      \item Use a generic matching method to obtain a matched sample (noting optimal), to ensure good balance
      \item Then use regression (e.g., bias-corrected) or weighting on the matched sample to correct for residual imbalance and improve efficiency
    \end{enumerate}
    \item Empirical results show that once a matched sample with reasonable balance is obtained, the point estimates from different additional methods are usually similar
    \item In high dimensional cases, combining matching (or stratification or weighting) with regression is the way to go (i.e., augmented or double learning)
    \item Matching is the starting point, not the end
  \end{itemize}

\end{frame}

\begin{frame}
  \frametitle{Mix Matching with Regression}

  \begin{itemize}
    \item Matching and regression are strongly connected, but why doing matching and regression separately?
    \item Residual imbalance in matching can be addressed via regression adjustment
    \item Rubin (1973): perform bias correction via regression on the matched sample
    \item Abadie and Imbens (2006): provided the theoretical basis and software for bias-corrected matching estimator
    \item Let $\mu_{d} (\mathbf{x}) = \mathbb{E} \left[ Y(d) \mid \mathbf{X}_i = \mathbf{x} \right]$, and $\hat{\mu}_d (\mathbf{X}_i)$ be consistent estimator of $\mu_d(\mathbf{X}_i)$, for $d = 0, 1$. A regression estimator uses $\hat{\mu}_d(\mathbf{X}_i)$ to impute missing potential outcomes $Y_i(d)$. 
  \end{itemize}

\end{frame}

\begin{frame}
  \frametitle{Bias-Corrected Matching: Formal Definition}

  \begin{itemize}
    \item Combine matching with regression adjustment for residual imbalance:
    \begin{align*}
      \tilde{Y}_i(0) = \begin{cases}
        Y_i, & D_i = 0, \\
        \frac{1}{M} \sum_{j \in \mathcal{M}_i} \left[ Y_j + \hat{\mu}_0(\mathbf{X}_i) - \hat{\mu}_0(\mathbf{X}_j) \right], & D_i = 1,
      \end{cases}\\
      \tilde{Y}_i(1) = \begin{cases}
        \frac{1}{M} \sum_{j \in \mathcal{M}_i} \left[ Y_j + \hat{\mu}_1(\mathbf{X}_i) - \hat{\mu}_1(\mathbf{X}_j) \right], & D_i = 0, \\
        Y_i, & D_i = 1.
      \end{cases}
    \end{align*}
    \item The bias-corrected matching estimators:
    \begin{align*}
      \hat{\tau}^{\text{bcm, ATE}} = \frac{1}{N} \sum_{i=1}^N \left( \tilde{Y}_i(1) - \tilde{Y}_i(0) \right)\\
      \hat{\tau}^{\text{bcm, ATT}} = \frac{1}{N_1} \sum_{i=1}^N \left( Y_i - \tilde{Y}_i(0) \right) D_i
    \end{align*}
  \end{itemize}

\end{frame}


\begin{frame}
  \frametitle{Connections between Matching and Regression}

  \begin{itemize}
    \item Both matching and regression estimators impute missing potential outcomes
    \item For unit $i$ with $D_i = 0$, we observed $Y_i(0)$ but need to impute $Y_i(1)$
    \item \textbf{Regression imputation}: uses $\hat{\mu}_1(\mathbf{X}_i)$ where $\hat{\mu}_d(\mathbf{x})$ is a consistent estimator of $\mu_d(\mathbf{x}) = \mathbb{E}\left[ Y \mid \mathbf{X}_i = \mathbf{x}, D = d \right]$
    \begin{align*}
      \tilde{Y}_i(1) = \begin{cases}
        \hat{\mu}_1(\mathbf{X}_i) & \text{if } D_i = 0, \\
        Y_i, & \text{if } D_i = 1
      \end{cases}
    \end{align*}
    \item \textbf{Matching}: uses average of $M$ nearest neighbors
    \begin{align*}
      \hat{Y}_i(1) = \begin{cases}
        \frac{1}{M} \sum_{j \in \mathcal{M}_i} Y_j & \text{if } D_i = 0, \\
        Y_i, & \text{if } D_i = 1
      \end{cases}
    \end{align*}
    \item Key insight: Matching is a nearest-neighbor estimate of $\mu_1(\mathbf{X}_i)$, but with fixed $M$ as $N \to \infty$, hence it does NOT consistently estimate $\mu_1(\mathbf{X}_i)$.
  \end{itemize}
    
\end{frame}

\begin{frame}
  \frametitle{From simple to Bias-Corrected Matching}

  \begin{itemize}
    \item Consider imputing $Y_i(1)$ for unit $i$ with $D_i = 0$
    \item Simple matching:
    \begin{align*}
      \hat{Y}_i(1) = \frac{1}{M} \sum_{j \in \mathcal{M}_i} Y_j
    \end{align*}
    \item Bias-corrected matching (Rubin 1973; Abadie and Imbens 2011):
    \begin{align*}
      \tilde{Y}_i(1) &= \frac{1}{M} \sum_{j \in \mathcal{M}_i} \left[ Y_j + \hat{\mu}_1(\mathbf{X}_i) - \hat{\mu}_1(\mathbf{X}_j) \right]\\
      &= \hat{Y}_i(1) + \underbrace{\left[ \hat{\mu}_1 (\mathbf{X}_i) - \frac{1}{M} \sum_{j \in \mathcal{M}_i} \hat{\mu}_1(\mathbf{X}_j) \right]}_{\text{bias correction}}
    \end{align*}
    \item Since $\norm{\mathbf{X}_i - \mathbf{X}_j}$ is small for $j \in \mathcal{M}_i$, adjustment is small
    \item But this small adjustment removes the bias of simple matching estimator asymptotically!
  \end{itemize}
  
\end{frame}


\begin{frame}
  \frametitle{An Alternative Representation via Bias-Corrected Regression}

  \begin{itemize}
    \item Bias-corrected matching can also be viewed from regression side:
    \item For unit $i$ with $D_i = 0$, the imputation is 
    \begin{align*}
      \tilde{Y}_i(1) = \underbrace{\hat{\mu}_1(\mathbf{X}_i)}_{\text{regression imputation}} + \underbrace{\frac{1}{M} \sum_{j \in \mathcal{M}_i} \left[ Y_j - \hat{\mu}_1(\mathbf{X}_j) \right]}_{\text{robustness via local averaging}}
    \end{align*}
    \item Interpretation:
    \begin{itemize}
      \item Start with regression estimator $\hat{\mu}_1(\mathbf{X}_i)$
      \item Add local correction using actual outcomes near $\mathbf{X}_i$
      \item The correction term provides robustness against misspecification
    \end{itemize}
  \end{itemize}
\end{frame}


\begin{frame}
  \frametitle{More Insights}

  \begin{itemize}
    \item Comparison of properties under misspecification
    \begin{itemize}
      \item Regression: if $\hat{\mu}_d$ is misspecified, estimator is inconsistent.
      \item Matching: robust but biased asymptotically
      \item Bias-corrected: robust against misspecification of regression function
    \end{itemize}
    \item The why (Abadie and Imbens 2011):
    \begin{itemize}
      \item Under weak assumptions, $\frac{1}{M} \sum_{j \in \mathcal{M}_i} Y_j \to \mu_d(\mathbf{X}_i)$
      \item If regression is misspecified, bias correction $\approx \mu_d(\mathbf{X}_i) - \hat{\mu}_d (\mathbf{X}_i)$
      \item This eliminates inconsistency from regression imputation
      \item If regression is correct, adding $\frac{1}{M} \sum_{j \in \mathcal{M}_i} \left[ Y_j - \hat{\mu}_d(\mathbf{X}_j) \right]$ only adds noise
    \end{itemize}
  \end{itemize}

\end{frame}


\begin{frame}
  \frametitle{Large Sample Properties of Matching Estimators}

  \begin{itemize}
    \item The simple matching estimator can be decomposed (Abadies and Imbens 2006):
    \begin{align*}
      \hat{\tau}_{M}^{\text{ATE}} - \tau = \underbrace{\left( \overline{\tau(\mathbf{X})} - \tau \right)}_{\text{sampling var.}} + \underbrace{D_M}_{\text{weighted residuals}} + \underbrace{B_M}_{\text{conditional bias}}
    \end{align*}
    where
    \begin{itemize}
      \item $\overline{\tau(\mathbf{X})} = \frac{1}{N} \sum_{i = 1}^N \left[ \mu_1(\mathbf{X}_i) - \mu_0(\mathbf{X}_i) \right]$ 
      \item $D_M = \frac{1}{N} \sum_{i = 1}^N \left( 2D_i - 1 \right) \left( 1 + \frac{K_M(i)}{M} \right)\varepsilon_i$, where $\varepsilon_i = Y_i - \mu_{D_i}(\mathbf{X}_i)$
      \item $B_M =\frac{1}{N} \sum_{i = 1}^N \frac{2D_i - 1}{M} \sum_{j = 1}^M \left[ \mu{1 - D_i}(\mathbf{X}_i) - \mu_{1 - D_i}(\mathbf{X}_{l_j(i)}) \right]$
      \item $K_M(i) = $ number of times unit $i$ is used as a match
    \end{itemize}
    \item First two terms: mean zero, $O_p \left( N^{-1 / 2} \right)$, asymptotically normal 
    \item Bias term $B_M$: $O_p \left( N^{-1 / p} \right)$ where $p = $ dimension of covariates
  \end{itemize}

\end{frame}

\begin{frame}
  \frametitle{Large Sample Properties of Matching Estimators}

  \begin{itemize}
    \item Conditional variance of simple matching (Abadie and Imbens 2006):
    \begin{align*}
      \Var \left( \hat{\tau}_{M}^{\text{ATE}} \mid \mathbf{X}, \mathbf{D} \right) = \Var \left( D_M \mid \mathbf{X}, \mathbf{D} \right)  = \frac{1}{N^2} \sum_{i = 1}^N \left( 1 + \frac{K_M(i)}{M} \right)^2 \sigma_{D_i}^2 (\mathbf{X_i})
    \end{align*}
    where $\sigma_{D_i}^2 (\mathbf{X_i}) = \Var \left( \varepsilon_i \mid \mathbf{X}_i = \mathbf{x}, D_i = d \right)$
    \item Let $V^E = N \Var \left( \hat{\tau}_{M}^{\text{ATE}} \mid \mathbf{X}, \mathbf{D} \right)$, and $V^{\tau(\mathbf{X})} = \mathbb{E}\left[ (\tau(\mathbf{X}_i) - \tau)^2 \right]$, where $\tau(\mathbf{x}) = \mathbb{E}\left[ Y_i(1) - Y_i(0) \mid \mathbf{X}_i = \mathbf{x} \right]$
    \item Asymptotic normality (Abadie and Imbens (2006), Theorem 1):
    \begin{align*}
      \left\{ V^E + V^{\tau(\mathbf{X})} \right\}^{-1/2} \sqrt{N} \left( \hat{\tau}_{M}^{\text{ATE}} - B_M - \tau \right) \xrightarrow{d} \mathcal{N}(0, 1)
    \end{align*}
    \item Crucial insight: simple matching itself only achieves $\sqrt{N}$-consistency when $p = 1$ (e.g., true propensity score) since $B_M = O_p \left( N^{-1 / p} \right)$. Bias-corrected matching removes $B_M$, achieving asymptotic normality for any $p$:
    \begin{align*}
      \left\{ V^E + V^{\tau(\mathbf{X})} \right\}^{-1/2} \sqrt{N} \left( \hat{\tau}_M^{\text{bcm, ATE}} - \tau \right) \xrightarrow{d} \mathcal{N}(0, 1)
    \end{align*}
  \end{itemize}

  

\end{frame}

\end{document}
